\chapter{Eserciziario per lo Scritto}

Questo capitolo contiene esercizi risolti passo-passo simili a quelli proposti nelle esercitazioni e all'esame scritto.

\section{Esercizio Tipo 1: Campionamento e Aliasing}
\textbf{Testo:} Dato il segnale $x(t) = \sin(2\pi \cdot 100 t)$, calcolare la frequenza di Nyquist.
Se il segnale viene campionato a $F_s = 150 \text{ Hz}$, qual è la frequenza alias percepita?

\vspace{0.5cm}
\textbf{Risoluzione:}

\begin{enumerate}
    \item Identificare la frequenza del segnale: Dalla formula $x(t) = \sin(2\pi f t)$, otteniamo $f_{max} = 100 \text{ Hz}$.
    \item Calcolare la frequenza di Nyquist: $F_{nyq} = 2 \cdot f_{max} = 2 \cdot 100 = 200 \text{ Hz}$.
    \item Calcolo dell'alias: Poiché $F_s = 150 \text{ Hz} < 200 \text{ Hz}$, non rispettiamo il teorema di Nyquist e si verificherà aliasing.
    La frequenza alias percepita si ottiene calcolando la distanza della frequenza del segnale dal multiplo più vicino di $F_s$:
    $f_{alias} = |F_s - f_{max}| = |150 - 100| = 50 \text{ Hz}$.
\end{enumerate}

\section{Esercizio Tipo 2: Quantizzazione, SNR e Bitrate}
\textbf{Testo:} Si vuole quantizzare un segnale con dinamica in ingresso tra 0 V e 5 V utilizzando un ADC a 8 bit.
Calcolare i livelli, lo step $\Delta$, l'approssimazione del rapporto segnale/rumore di quantizzazione (SNR) in dB, e il bitrate se campionato a $F_s = 200 \text{ Hz}$.
Cosa succede se si passa da 8 a 16 bit?

\vspace{0.5cm}
\textbf{Risoluzione:}

\begin{enumerate}
    \item \textbf{Livelli ($L$):} $L = 2^b = 2^8 = 256$ livelli.
    \item \textbf{Step $\Delta$:} $\Delta = \frac{V_{max} - V_{min}}{L - 1} = \frac{5 - 0}{255} \approx 0.0196 \text{ V}$.
    \item \textbf{SNR (SQNR):} $SQNR_{dB} \approx 6.02 \cdot b + 1.76 = 6.02 \cdot 8 + 1.76 = 49.92 \text{ dB}$.
    \item \textbf{Bitrate:} Il bitrate si definisce come $R = b \cdot F_s = 8 \cdot 200 = 1600 \text{ bps}$ (bit per secondo).
    \item \textbf{Passaggio a 16 bit:} L'SNR diventa $6.02 \cdot 16 + 1.76 = 98.08 \text{ dB}$, con una notevole diminuzione del rumore (incremento di qualità).
    Il bitrate raddoppia ($16 \cdot 200 = 3200 \text{ bps}$).
\end{enumerate}

\section{Esercizio Tipo 3: Filtraggio 1D (Array)}
\textbf{Testo:} Dato l'array $x = [2, 3, 2, 20, 3, 2, 1]$, applicare un filtro Mediano (finestra $W=3$) e un filtro di Media Mobile (Media, $W=3$, con zero padding o mantenendo i bordi inalterati).
Spiegare le differenze.

\vspace{0.5cm}
\textbf{Risoluzione:}

Consideriamo i punti interni per semplicità.
Valutiamo $i=3$ (valore 20) che è chiaramente uno spike.

\begin{itemize}
    \item \textbf{Filtro Mediano ($W=3$):}
    \begin{itemize}
        \item Centrato sul 20 ($x[3], x[4], x[5]$): $\{2, 20, 3\} \rightarrow \text{ordinato } \{2, 3, 20\} \rightarrow \text{Mediano} = 3$.
        \item Centrato sul secondo 2 ($x[2], x[3], x[4]$): $\{3, 2, 20\} \rightarrow \text{ordinato } \{2, 3, 20\} \rightarrow \text{Mediano} = 3$.
    \end{itemize}
    Il valore $y$ mediano scarta l'outlier "20" considerandolo un rumore impulsivo. $y \approx[2, 3, 3, 3, 2, 2, 1]$.

    \item \textbf{Filtro Media ($W=3$):}
    \begin{itemize}
        \item Centrato sul 20: $\frac{2 + 20 + 3}{3} = \frac{25}{3} \approx 8.33$.
        \item Centrato sul secondo 2: $\frac{3 + 2 + 20}{3} = \frac{25}{3} \approx 8.33$.
    \end{itemize}
    Il picco si è schiacciato ($\sim 8.3$ invece di $20$), ma si è "spalmato" sui campioni vicini alterandoli nettamente.
\end{itemize}

\section{Esercizio Tipo 4: Valutazione Modelli (Matrice di Confusione)}
\textbf{Testo:} Dato il risultato di un modello di rilevamento Spam (o Malattia), con $TP = 40, TN = 950, FP = 10, FN = 5$.
Calcolare Accuracy, Precision, Recall e F1-Score.
Spiegare quando l'Accuracy è ingannevole.

\vspace{0.5cm}
\textbf{Risoluzione:}

\begin{itemize}
    \item \textbf{Accuracy:} $\frac{TP + TN}{TP + TN + FP + FN} = \frac{40 + 950}{1005} = \frac{990}{1005} = 0.985$ (98.5\%).
    \item \textbf{Precision:} $\frac{TP}{TP + FP} = \frac{40}{40 + 10} = \frac{40}{50} = 0.80$ (80\%).
    \item \textbf{Recall (Sensitivity):} $\frac{TP}{TP + FN} = \frac{40}{40 + 5} = \frac{40}{45} \approx 0.888$ (88.8\%).
    \item \textbf{F1-Score:} $2 \cdot \frac{Precision \cdot Recall}{Precision + Recall} \approx 2 \cdot \frac{0.8 \cdot 0.888}{0.8 + 0.888} \approx 0.84$.
\end{itemize}

\textbf{Analisi:} L'Accuracy è molto alta (98.5\%) e ingannevole.
Questo perché il dataset è estremamente sbilanciato (Class Imbalance): i veri negativi ($TN$) dominano.
Se un classificatore dicesse sempre "Non Spam", avrebbe un accuracy di $\frac{960}{1005} = 95.5\%$, pur non avendo capacità predittiva.
In questi casi, F1-Score e Recall sono metriche più significative.

\section{Esercizio Tipo 5: Energia e Potenza di una Sinusoide}
\textbf{Testo:} Un segnale $x(t) = 10 \cos(2\pi f t)$ viene osservato su un singolo periodo $T$.
Qual è l'energia $E$ contenuta in quel periodo?
Se la frequenza raddoppia, cosa succede alla potenza media su tempo infinito?

\vspace{0.5cm}
\textbf{Risoluzione:}
\begin{enumerate}
    \item L'energia su un singolo periodo per una sinusoide di ampiezza $A$ è data dall'integrale del quadrato, che equivale a $E_{periodo} = \frac{A^2}{2} \cdot T$. Sostituendo $A=10$, otteniamo $E = \frac{100}{2}T = 50 T$.
    \item La potenza media teorica di una sinusoide su tempo infinito dipende \textit{solo} dall'ampiezza ed è $P = \frac{A^2}{2}$.
    Essendo indipendente da $f$, se la frequenza raddoppia, \textbf{la potenza rimane invariata}.
\end{enumerate}

\section{Esercizio Tipo 6: Filtraggio di Segnali Fisiologici}
\textbf{Testo:} Stai analizzando un segnale EDA/GSR. Il soggetto tocca accidentalmente l'elettrodo per un millisecondo, creando un picco stretto e altissimo (artefatto impulsivo).
Vuoi estrarre la linea di base (SCL) senza distorcerla.
Hai a disposizione un filtro di media mobile e un filtro mediano.
Qual è la strategia/sequenza corretta?

\vspace{0.5cm}
\textbf{Risoluzione:}
La sequenza corretta è: \textbf{Prima Filtro Mediano, poi Filtro di Media (o estrazione del trend)}.
\textit{Motivazione:} L'artefatto impulsivo (outlier) verrebbe "spalmato" da un filtro di media, inquinando temporalmente il segnale SCL.
Applicando prima un filtro non-lineare (Mediano), il picco anomalo viene scartato completamente.
A quel punto, sul segnale pulito, si può applicare un filtro passa-basso (media) per estrarre la linea di base.

\section{Esercizio Tipo 7: Domande Teoriche di Machine Learning}

\textbf{Domanda 1:} In un classificatore Random Forest, perché le performance sono solitamente migliori di un singolo Decision Tree?

\vspace{0.25cm}
\textit{Risposta:} Perché essendo un algoritmo \textit{Ensemble}, combina molti alberi diversi addestrati su sottoinsiemi randomici di dati e feature.
Questo approccio basato sul voto a maggioranza \textbf{riduce la varianza} complessiva del modello e previene l'Overfitting tipico dei singoli alberi decisionali, che tendono ad imparare il rumore.

\vspace{0.5cm}
\textbf{Domanda 2:} Hai un dataset a forma di cerchio: i punti della classe A sono al centro, i punti della classe B formano un anello esterno.
Come procedi con una SVM?

\vspace{0.25cm}
\textit{Risposta:} Un iperpiano lineare non può separare le classi.
Devo usare un \textbf{Kernel RBF (Gaussiano)} o Polinomiale.
Il kernel proietta (trasforma) i dati in un nuovo spazio matematico a più dimensioni, rendendo le due distribuzioni linearmente separabili in quella nuova dimensionalità.

\vspace{0.5cm}
\textbf{Domanda 3:} Cosa indica il fenomeno dell'Overfitting in termini di Bias e Varianza?

\vspace{0.25cm}
\textit{Risposta:} L'Overfitting si manifesta con \textbf{Basso Bias e Alta Varianza}.
Il modello è talmente complesso che cattura perfettamente i dati di training (incluso il loro rumore casuale), ma diventa ipersensibile alle fluttuazioni, fallendo sui nuovi dati di test.
